%%%%%%%%%%  scan idx 350  |  folio 337  %%%%%%%%%%
\SGAPrintPage{337}%
generating the canonical filtration.

Writing \(b=\sum a_i m_i\), with \(a_i\in A\) and the \(m_i\) monomials in the \(\gamma^j X\), and using the additivity of \(\gamma_t\), we see that it suffices to give the proof for \(\gamma^k(am_i)\), with \(a\in A\) and \(m_i\in M_i\).

Using Lemma 4.11. and the \(A\)-\(\lambda\)-homomorphism from \(B\) to \(B\) which sends \(X\) to \(m_i\), we see that it suffices to give the proof for an element of the form \(\gamma^k(m_i)\).

Finally, since the polynomial \(P'_k\) has no isobaric components of weight \(<k\), we are reduced to the case \(\gamma^k(\gamma^j X)\), which itself follows from the fact that the polynomials \(Q'_{j,k}\) have no isobaric components of weight \(<k\) (3.6).

\sgaStableTarget{sga6.e05.rem.4-13}\textbf{Remarks 4.13.} i) The decomposition of \(\Fil^k B\) given in 4.12 is a direct-sum decomposition of \(A\)-modules. In other words, if \(b\in\Fil^k B\), and if we write \(b=\sum_{j\le k}a_{i,j}m_{i,j}+b'\), with \(a_{i,j}\in A\), \(m_{i,j}\in M_j\), and \(b'\in B\mathbin{.}M_{k+1}\), then \(a_{i,j}\in\Fil^{k-j}A\).

ii) Statement 4.12. remains true for every quotient of \(B\) by a \(\lambda\)-ideal contained in the augmentation ideal.

iii) The statement analogous to 4.12., obtained by replacing the indeterminate \(X\) by a family of indeterminates \((X_i)\), is also true.

\sgaStableTarget{sga6.e05.cor.4-14}\textbf{Corollary 4.14.} \sgaStableTarget{sga6.e05.eq.4-14-1}\(\Gr^k B\approx\Gr^k A\oplus\Gr^{k-1}A\mathbin{.}M_1\oplus\ldots\oplus\Gr^0 A\mathbin{.}M_k\).

We define an arrow from right to left by associating to each element \(\bar a\mathbin{.}m\), where \(\bar a\) is the class in \(\Gr^{k-i}A\) of an element \(a\) of \(\Fil^{k-i}A\) and \(m\in M_i\), the class in \(\Gr^k B\) of \(am\). We define an arrow in the opposite direction by decomposition according to Proposition 4.12 and Remark 4.13.i). It is clear that these two arrows are inverse to one another.

%%%%%%%%%%  scan idx 351  |  folio 338  %%%%%%%%%%
\SGAPrintPage{338}%
\sgaStableTarget{sga6.e05.cor.4-15}\textbf{Corollary 4.15.} \uline{Let} \(B'\) \uline{be the} \(\lambda\)-\uline{ring generated over} \(K\) \uline{by the indeterminate} \(X\) \uline{subject to the relations} \(\gamma^k(X)=0\) \uline{for} \(k>n\), \uline{i}.\uline{e}. \(B'=K[X,\ldots,\gamma^n X]\) (4.9). \uline{Then the} \(\lambda\)-\uline{ring filtration of} \(B'\) \uline{is identical to the weight filtration with respect to the} \(\gamma^j X\).

Taking account of Proposition 4.12, Remark 4.13.ii), and the fact that \(A=K\), we have \(\Fil^k B'=K\mathbin{.} M_k+B'\mathbin{.} M_{k+1}\), so \(\Fil^k B'\) is the ideal generated by the monomials of weight \(\ge k\) with respect to the \(\gamma^j X\).

\sgaStableTarget{sga6.e05.sec.5}\section*{5. The \(\lambda^p(N,x)\).}

This section will be used only in Exp.VII~4, for the study of regular immersions; it may therefore be omitted on first reading.

\sgaStableTarget{sga6.e05.para.5-1}5.1. Let \(d\) be an integer \(\ge1\). We denote by \(A\) the \(\mathbb Z\)-\(\lambda\)-algebra generated by the generators \(N\), \(x\) subject to the relations \(\lambda^i(N)=0\) for \(i>d\); we therefore have:
\[
  A=\mathbb Z[N,\ldots,\lambda^dN;
       x,\ldots,\lambda^jx,\ldots].
\]

We put \sgaStableTarget{sga6.e05.eq.5-1-1}\(\lambda_{-1}(N)=\sum_{i=0}^{d}(-1)^i\lambda^i(N)\).

\sgaStableTarget{sga6.e05.lem.5-2}\textbf{Lemma 5.2.} \uline{For every} \(p\ge1\), \(\lambda^p(x\mathbin{.}\lambda_{-1}(N))\) \uline{is divisible by} \(\lambda_{-1}(N)\).

We know (2.3) that \(\lambda^p(x\mathbin{.}\lambda_{-1}(N))=P_p(\lambda^i x,\lambda^j(\lambda_{-1}(N)))\) is an isobaric polynomial of weight \(p\) with respect to the \(\lambda^j(\lambda_{-1}(N))\); it therefore suffices to prove the lemma for the \(\lambda^p(\lambda_{-1}(N))\).

Let \(B\) be the \(\lambda\)-ring generated over \(\mathbb Z\) by the indeterminates \(N_1,\ldots,N_d\), subject to the relations \(\lambda^i(N_j)=0\) if \(i\ge2\). It is clear that \(\lambda^k(\sum_{i=1}^{d}N_i)=0\) for \(k>d\). We can therefore define a \(\lambda\)-homomorphism \(f\) from \(A\) to \(B\) by sending \(N\) to \(\sum_{i=1}^{d}N_i\); we then have \(\Ghat(f)(\lambda_t(N))=\prod_{i=1}^{d}\lambda_t(N_i)=\prod_{i=1}^{d}(1+N_it)\),
%%%%%%%%%%  scan idx 352  |  folio 339  %%%%%%%%%%
\SGAPrintPage{339}%
which shows that, for \(1\le k\le d\), \(\lambda^k N\) maps to the \(k\)-th symmetric function of the \(N_i\), and consequently that \(f\) is injective and identifies \(A\) with the ring of symmetric polynomials in the \(N_i\).

We then have \(\lambda_{-1}(N)=\prod_{i=1}^{d}(1-N_i)\), and \(\lambda_t(\lambda_{-1}(N))=\lambda_t\bigl(\prod_{i=1}^{d}(1-N_i)\bigr)\). The \(\lambda^p(\lambda_{-1}(N))\) are polynomials in the \(N_i\); if, for some index \(i\), we set \(N_i=1\), we see that \(\lambda_t(\lambda_{-1}(N))\) becomes \(\lambda_t(0)=1\), hence that the \(\lambda^p(\lambda_{-1}(N))\) vanish for \(N_i=1\), and therefore that they contain \(1-N_i\) as a factor. Thus they contain the product of the \(1-N_i\) as a factor, i.e. they are divisible by \(\lambda_{-1}(N)\) in \(B\), and hence also in \(A\), since the quotient of two symmetric polynomials is evidently symmetric.

\sgaStableTarget{sga6.e05.def.5-3}5.3. Since \(A\) is an integral domain, we define an element \(\lambda^p(N,x)\) unambiguously by putting (\(p\ge1\)):
\sgaStableTarget{sga6.e05.eq.5-3-1}\begin{equation}\tag{5.3.1}
  \lambda^p(N,x)\mathbin{.}\lambda_{-1}(N)=\lambda^p(x\mathbin{.}\lambda_{-1}(N)).
\end{equation}

The element \(\lambda^p(N,x)\) is therefore a polynomial in the \(\lambda^i(x)\), \(\lambda^j(N)\), isobaric of weight \(p\) with respect to the \(\lambda^i(x)\), since the same is true of \(\lambda^p(x\mathbin{.}\lambda_{-1}(N))\) by 2.3.

Let \(d'\) be an integer \(\ge d\), and let \(A'\) be the \(\mathbb Z\)-\(\lambda\)-algebra generated by the generators \(N'\) and \(x'\) subject to the relations \(\lambda^i(N')=0\) for \(i>d'\); by Proposition 4.7, there exists exactly one \(\lambda\)-homomorphism \(f:A'\longrightarrow A\), sending \(N'\) to \(N\) and \(x'\) to \(x\). Definition (5.3.1) implies that \(f(\lambda^p(N',x'))=\lambda^p(N,x)\) for every \(p\ge1\).

If \(B\) is a \(\lambda\)-ring and \(N''\) is an element of \(B\) such that \(\lambda^k(N'')=0\) for \(k>d\), we can, for every \(x''\in B\), define an element \(\lambda^p(N'',x'')\) by substitution, or by using the unique \(\lambda\)-homomorphism from \(A\) to
%%%%%%%%%%  scan idx 353  |  folio 340  %%%%%%%%%%
\SGAPrintPage{340}%
\(B\) which sends \(N\) to \(N''\), and \(x\) to \(x''\); by the preceding remark, the element obtained does not depend on the integer \(d\) such that \(\lambda^k(N'')=0\) for \(k>d\). Using the same homomorphism, we see that relation (5.3.1) is satisfied by \(N''\) and \(x''\).

\sgaStableTarget{sga6.e05.para.5-4}5.4. Consider the \(\lambda\)-ring generated over \(\mathbb Z\) by the indeterminates \(N,x,y\), subject to the relations \(\lambda^k(N)=0\) for \(k>d\). From (5.3.1) we obtain:
\begin{align*}
  \lambda^p(N,x+y)\mathbin{.}\lambda_{-1}(N)
  &=\lambda^p\bigl((x+y)\mathbin{.}\lambda_{-1}(N)\bigr)\\
  &=\sum_{i+j=p}\lambda^i(x\mathbin{.}\lambda_{-1}(N))\mathbin{.}\lambda^j(y\mathbin{.}\lambda_{-1}(N))\\
  &=\sum_{\substack{i+j=p\\ i,j\ne0}}\lambda^i(N,x)\mathbin{.}\lambda^j(N,y)\mathbin{.}(\lambda_{-1}(N))^2
    +(\lambda^p(N,x)+\lambda^p(N,y))\lambda_{-1}N
\end{align*}

Since the ring under consideration is an integral domain, we deduce the following relation, valid by substitution in every \(\lambda\)-ring, and for every element \(N\) such that \(\lambda^k(N)=0\) for \(k>d\):
\sgaStableTarget{sga6.e05.eq.5-4-1}\begin{equation}\tag{5.4.1}
  \lambda^p(N,x+y)=\sum_{\substack{i+j=p\\ i,j\ne0}}\lambda^i(N,x)\mathbin{.}\lambda^j(N,y)\mathbin{.}\lambda_{-1}(N)+\lambda^p(N,x)+\lambda^p(N,y).
\end{equation}

Now consider the \(\lambda\)-ring generated over \(\mathbb Z\) by indeterminates \(N,N',x\) subject to the relations \(\lambda^k(N)=0\) for \(k>d\), and \(\lambda^k(N')=0\) for \(k>d'\). We then have \(\lambda^k(N+N')=0\) for \(k>d+d'\); we can therefore define the \(\lambda^p(N+N',x)\). We obtain:
\[
  \lambda^p(N+N',x)\mathbin{.}\lambda_{-1}(N+N')=\lambda^p(x\mathbin{.}\lambda_{-1}(N+N')).
\]
By the additivity of \(\lambda_t\), we have: \(\lambda_{-1}(N+N')=\lambda_{-1}(N)\mathbin{.}\lambda_{-1}(N')\). We deduce:
\begin{align*}
  \lambda^p(N+N',x)\mathbin{.}\lambda_{-1}(N)\mathbin{.}\lambda_{-1}(N')
  &=\lambda^p(x\mathbin{.}\lambda_{-1}(N)\mathbin{.}\lambda_{-1}(N'))\\
  &=\lambda^p(N',x\mathbin{.}\lambda_{-1}(N))\mathbin{.}\lambda_{-1}(N').
\end{align*}

%%%%%%%%%%  scan idx 354  |  folio 341  %%%%%%%%%%
\SGAPrintPage{341}%
Since the ring under consideration is an integral domain, we deduce the following relation, valid by substitution in every \(\lambda\)-ring, and for all elements \(N,N'\) such that \(\lambda^k(N)=0\) for \(k>d\), \(\lambda^k(N')=0\) for \(k>d'\);
\sgaStableTarget{sga6.e05.eq.5-4-2}\begin{equation}\tag{5.4.2}
  \lambda^p(N+N',x)\mathbin{.}\lambda_{-1}(N)=\lambda^p(N',x\mathbin{.}\lambda_{-1}(N)).
\end{equation}

\sgaStableTarget{sga6.e05.para.5-5}5.5. Now let \(A\) be any \(\lambda\)-ring, and let \(N\) be an element of \(A\) such that \(\lambda^k(N)=0\) for \(k>d\). We shall associate to \(A\) a new \(\lambda\)-ring \(A_N\).

We first define a new multiplication on \(A\) by putting, for \(x,y\in A\):
\begin{equation}\tag{5.5.1}
  x\mathbin{.}_N y=x\mathbin{.} y\mathbin{.}\lambda_{-1}(N).
\end{equation}

The ring \(A\) with this multiplication is no longer unital; we therefore consider the ring \(\mathbb Z\times A\), obtained by adjoining a unit to \(A\). Its operations are thus defined by:
\begin{equation}\tag{5.5.2}
\begin{aligned}
  (n,x)+(n',x')&=(n+n',x+x')\\
  (n,x)\mathbin{.}(n',x')&=(nn',nx'+n'x+x\mathbin{.}_N x')
\end{aligned}
\end{equation}
We regard \(A\) as embedded in this ring by \(x\rightsquigarrow(0,x)\). We then define a pre-\(\lambda\)-ring structure by:
\sgaStableTarget{sga6.e05.eq.5-5-1}\begin{equation}\tag{5.5.3}
  \lambda_t(n,x)=(1+t)^n\mathbin{.}\sum_{p=0}^{\infty}\lambda^p(N,x)\mathbin{.} t^p.
\end{equation}
This is a pre-\(\lambda\)-ring structure by (5.4.1). We shall denote by \(A_N\) the pre-\(\lambda\)-ring thus obtained. To see that it is a \(\lambda\)-ring, it suffices by substitution to give the proof in the case where \(A\) is the \(\lambda\)-ring generated over \(\mathbb Z\) by the generators \(N,X,Y\), subject to the relations \(\lambda^k(N)=0\) for \(k>d\). We then define a homomorphism \(A_N\to A\) by
